\documentclass[numbered]{trbunofficial}
\usepackage{amssymb}
\usepackage{booktabs}
\usepackage{graphicx}
\usepackage{lastpage}
\usepackage[hidelinks]{hyperref}

\title{Latency-Aware Truckload Bid Acceptance under Operational
Feasibility: A Public Benchmark with Hindsight Ceilings}
\AuthorHeaders{Chandrasekaran}
\author{%
Aswin Chandrasekaran*\\
Bubba AI\\
\texttt{<FILL: City, State/Country, Postcode>}\\
ORCID: \texttt{<FILL if available>}\\
Email: aswin@bubba.ai\\[6pt]
*Corresponding Author%
}
\date{August 1, 2026}

% Current TRB sample lists total pages on the title page (not a word count).
% Set textwords to suppress the class's texcount call (needs no shell-escape),
% then override the title-page line to report total pages.
\setcounter{textwords}{2700}
\renewcommand{\printwordcount}{Total Number of Pages: \pageref{LastPage}}
\setcitestyle{aysep={}}  % TRB author-date style: (Author 2020), no comma

\begin{document}
\maketitle

\section*{Abstract}
\textbf{Objectives:} Online truckload bid acceptance requires accept or
reject decisions within seconds under closed-loop fleet state, yet public,
reproducible benchmarks are scarce and policies are usually judged only
against a suboptimal heuristic. We introduce a public benchmark and
measured performance ceilings for this problem and evaluate a
latency-aware acceptance policy.
\textbf{Methods:} We formulate a finite-horizon closed-loop Markov decision
process with operational feasibility (pickup reach, appointment windows,
simplified hours-of-service, stochastic yard delays) and three economic
layers (service-failure penalty, terminal fleet value, price-premium
window), calibrated from public Freight Analysis Framework (FAF) and
U.S. Department of Agriculture (USDA) truck-rate data. We derive two
hindsight ceilings---an exact small-prefix dynamic program and a
Lagrangian-per-truck information-relaxation bound---and evaluate a
surrogate-to-rollout cascade across ten paired seeds with bootstrap
confidence intervals.
\textbf{Findings:} The Lagrangian bound is 20.7\% tighter than an LP
relaxation on a tight-capacity scenario and 39.3\% tighter on a
scarce-capacity scenario, relocating the rollout teacher to 66--68\% of a
structure-respecting ceiling. The cascade recovers roughly 98\% of rollout
profit at 40--56\% of its mean decision latency and is statistically
indistinguishable from rollout on the tight scenario.
\textbf{Novelty:} This is a public-calibrated, dependency-free, closed-loop
benchmark that makes operational feasibility part of the reward and pairs
it with a structure-respecting hindsight ceiling.
\textbf{Practical Applications:} Feasibility-blind tender scoring carries a
hidden linear cost; a cheap model should be deployed with selective
escalation of high-stakes capacity decisions; and the remaining headroom
lies in cross-truck coordination rather than single-tender scoring.

\vspace{12pt}
\noindent\textit{Keywords}: truckload bid acceptance, closed-loop
stochastic optimization, approximate dynamic programming,
information-relaxation bounds, latency-aware policy cascade, public
benchmark.

\newpage

\section{Introduction}
Truckload carriers and brokers face an online accept/reject decision each
time a load tender arrives, and in a typical broker workflow the tender
must be priced and answered within seconds. The decision is operationally
constrained: a load can be served only by a truck that can reach the
pickup, satisfies appointment windows and hours-of-service (HOS) clocks,
and remains feasible through delivery. It is also economically
constrained: a tender with positive immediate margin can still be a poor
decision if it consumes a scarce truck before a predictable price-premium
window or strands the fleet in a market with weak outbound flow.
Conversely, a benchmark cannot claim to compare future-aware acceptance
policies meaningfully unless it (i) makes operational feasibility part of
the reward function rather than a side diagnostic, and (ii) anchors policy
evaluation against a measured ceiling rather than only against a finite
stochastic teacher.

Two gaps have held the field back. First, no public, reproducible
benchmark captures the closed-loop, fleet-state-dependent, latency-aware
truckload accept/reject problem: static routing instances are not
closed-loop, and dynamic-fleet studies typically rely on proprietary
operator data. Second, policy studies commonly report retention against a
rollout or reoptimization heuristic, which is itself suboptimal, leaving
unclear how much value remains relative to the true optimum.

This paper addresses both. Our contributions are: (1) a
public-calibrated, dependency-free, closed-loop benchmark
(FreightBidBench) with explicit operational feasibility and three economic
layers, reproducible from public Freight Analysis Framework (FAF) and
U.S. Department of Agriculture (USDA) truck-rate data; (2) two hindsight
ceilings---an exact small-prefix dynamic program and a Lagrangian-per-truck
information-relaxation bound substantially tighter than a naive LP
relaxation; and (3) a latency-aware surrogate-to-rollout cascade that
recovers nearly all of rollout's profit at a fraction of its decision
latency. Throughout, the benchmark is the primary artifact; the methods
are evidence that the sharpened problem admits a non-trivial
latency--profit frontier and a reproducible test bed for future work.

\section{Related Work}
Dynamic load acceptance in truckload operations has been studied as an
online state-dependent control problem, including large-fleet acceptance
with time windows \citep{kim2004dynamic} and real-time reoptimization
\citep{yang2004realtime}. The approximate dynamic programming (ADP) line
for fleet management \citep{simao2009adp,powell2007stochastic} established
that state-dependent value-function approximation can match operator-scale
decisions, but this work has historically been driven by private operator
data. Static vehicle-routing benchmarks \citep{solomon1987vrptw,uchoa2017cvrp}
shaped the routing literature yet are neither closed-loop nor stochastic in
tender arrival, and dynamic-VRP methods
\citep{ulmer2019offline,pillac2013review} target multi-customer routing
rather than online bid acceptance. Anchoring policy evaluation against
upper bounds via information relaxation \citep{brown2010information} and
Lagrangian relaxation of weakly coupled dynamic programs
\citep{adelman2008relaxations,yezhu2014weakly} is well established in
stochastic dynamic programming; we adopt this framing to bound the
truckload problem. Finally, cascading a cheap model to an expensive one,
escalating only uncertain instances, descends from cascade classifiers and
anytime/metareasoning control \citep{violajones2001,zilberstein1996}; what
has been missing is a systematic, publicly calibrated, feasibility-aware
instantiation for online freight decisions.

\section{Problem and Benchmark}
We model closed-loop truckload bid acceptance as a finite-horizon Markov
decision process over continuous, time-stamped events.

\subsection{State, Action, and Feasibility}
The horizon is $T=72$ hours and the fleet has $K$ trucks. Each truck's
state is $(\ell, \tau, h, d)$: its market, next-available time, and
remaining HOS drive and duty budgets. A load tender presents origin and
destination markets, posted price, linehaul cost, mileage, pickup and
delivery appointment windows, and stochastic yard-delay draws. The system
state is the fleet, the current tender, and time; the action is binary
(accept or reject). A deterministic feasibility map attempts to assign an
accepted load to a truck in the origin market, dropping trucks that miss
the latest pickup, computing pickup reach, HOS-feasible schedules with
rest insertions, and delivery arrival; among feasible trucks it selects the
one that becomes available earliest after delivery. If none is feasible the
fleet is not mutated and the accept is recorded as a service failure.

\subsection{Economic Layers}
Three versioned reward layers sharpen the decision problem so that distinct
policy classes separate. (L1) A \textit{service-failure penalty} $\rho$
(frozen at \$10) charges any accept of an infeasible load, so
feasibility-blind policies pay regret linear in their infeasible-accept
rate. (L2) A \textit{terminal fleet value} $\omega \sum_k V(\ell^{(k)}_T)$
with $\omega=0.25$ rewards ending trucks in high-value markets, where
$V(\cdot)$ combines FAF outbound intensity and a FAF/USDA net-outbound
imbalance panel; it penalizes greedy positioning. (L3) A daily
$1.5\times$ \textit{price-premium window} penalizes future-blind timing,
since accepting a moderate off-peak load can consume capacity needed for a
predictable on-peak premium.

\subsection{Public Calibration}
Candidate loads are drawn proportional to lane FAF tonnage; the initial
fleet is placed proportional to per-origin FAF outbound tonnage; distances
follow FAF ton-miles; and posted prices are drawn from USDA Agricultural
Marketing Service truck-rate bands. Implied per-mile rates (median \$3.28,
interquartile range \$2.83--\$3.57) match published refrigerated truckload
spot rates. The runtime uses only the Python standard library, and every
result is reproducible from a manifest recording the command, version
strings, seed pairs, and input paths.

\section{Hindsight Ceilings}
Retention against a finite rollout teacher does not reveal how much profit
is left relative to the optimal policy value $V^\star$, which is
intractable to compute for a fleet of dozens of trucks over a continuous
horizon. We therefore bound $V^\star$ from above.

\subsection{Exact Small-Prefix Dynamic Program}
For a truncated realized stream of $L$ loads, the exact problem is a binary
decision tree over deterministic post-acceptance fleet transitions, solved
by memoized search keyed on (prefix index, fleet hash), tractable for
$L \leq 16$. On a 12-load tight prefix the exact ceiling is matched by
simple feasibility-aware policies, confirming short prefixes are flat under
the v0.3 economics; the exact program is a trustworthy small-instance
reference rather than a full-horizon ceiling.

\subsection{Lagrangian-per-truck Bound}
The only constraint coupling trucks is that at most one truck serves each
load. Revealing the realized scenario and dropping integrality, sequencing,
and location constraints yields a valid but loose LP-style bound---an
information relaxation \citep{brown2010information}. Dualizing only the
assignment constraint with multipliers $\lambda_t \geq 0$ decouples the
fleet into independent per-truck sub-MDPs that retain each truck's HOS,
location, and sequencing structure:
\begin{equation}
L(\boldsymbol\lambda) = \sum_t \lambda_t + \sum_k V^k_{\boldsymbol\lambda}(u^{(k)}_0),
\qquad V^\star \leq \min_{\boldsymbol\lambda \geq 0} L(\boldsymbol\lambda).
\end{equation}
Weak duality gives validity for any $\boldsymbol\lambda \geq 0$; the joint
supremum decomposes across trucks because rewards and actions are
truck-separable once assignment is dualized. Each per-truck sub-MDP is
solved exactly by bucket-snapped forward enumeration---clocks are rounded
to each bucket's favorable corner, so per-truck values only over-estimate
and the joint upper bound is preserved---and $\boldsymbol\lambda$ is
optimized by diminishing-step subgradient ascent. Full proofs and the
looseness decomposition appear in the extended version. Empirically the
Lagrangian bound is 20.7\% tighter than the LP bound on \texttt{tight} and
39.3\% tighter on \texttt{scarce} (Table~\ref{tab:bound}).

\section{Policies: Surrogate and Cascade}
The policy set includes simple baselines (reject-all, accept-all-feasible,
myopic-margin, bid-price); a dependency-free linear \textit{surrogate} fit
by closed-form ridge regression on rollout incremental-value labels, using
standardized features and a feasibility guard; a finite \textit{rollout
teacher} that expands accept and reject branches under common random
numbers; and a \textit{cascade}. The cascade runs the cheap surrogate by
default and escalates to the rollout teacher only when a decision is
uncertain or high-stakes. Let $\Delta_\theta(s)$ be the surrogate's signed
accept-minus-reject score and $n_o(s)$ the number of immediately available
trucks in the load's origin market. With boundary band $\beta \geq 0$ and
scarcity threshold $\kappa$, escalation fires when
\begin{equation}
E(s;\beta,\kappa) = \mathbf{1}\{|\Delta_\theta(s)| \leq \beta\} \;\vee\; \mathbf{1}\{n_o(s) \leq \kappa\} = 1.
\end{equation}
The rollout-call share, and hence mean decision latency, is monotone in
both thresholds, so $(\beta,\kappa)$ traces a latency--profit frontier. The
scarcity trigger defers exactly the high-stakes capacity decisions on which
the surrogate is least reliable.

\section{Experiments}
All experiments use the frozen \texttt{scenario-v0.3.2} contract with ten
paired (train, eval) seeds; confidence intervals are paired-bootstrap 95\%
intervals from 20{,}000 resamples. We report a tight-capacity and a
scarce-capacity scenario, with a mild scenario as a negative control.
Layer-ablation sweeps confirm each layer separates its target policy class:
under $\rho=\$10$, feasibility-blind policies fall below the
feasibility-aware greedy baseline; $\omega=0.25$ reduces
accept-all-feasible retention below 95\% of rollout; and the price-premium
window widens the rollout-versus-best-simple gap past ten percentage
points.

\begin{table}[t]
\caption{Ten-Seed Methods Comparison under scenario-v0.3.2. Retention is
mean closed-loop profit relative to the rollout teacher; the cascade uses
band $\beta=\$500$ at scarcity threshold $\kappa=2$. The last column is the
paired-bootstrap 95\% CI on the cascade$-$rollout profit delta ($n=10$); an
interval containing \$0 means the cascade is statistically indistinguishable
from rollout.}
\label{tab:methods}
\centering\small
\resizebox{\textwidth}{!}{%
\begin{tabular}{lrrrrr}
\toprule
Scenario & Best simple & Surrogate & Cascade & Cascade/Rollout ms & $\Delta$ vs rollout (95\% CI) \\
\midrule
\texttt{tight}  & 91.0\% & 94.2\% & 98.2\% & 12.95 / 32.11 & $-\$23.3$k $[-\$61.2$k, $+\$11.6$k$]$ \\
\texttt{scarce} & 86.5\% & 89.3\% & 98.0\% & 11.54 / 20.59 & $-\$20.7$k $[-\$33.4$k, $-\$7.2$k$]$ \\
\texttt{mild}   & 100.1\% & 98.2\% & 99.7\% &  8.31 / 49.50 & $-\$4.2$k $[-\$18.2$k, $+\$11.0$k$]$ \\
\bottomrule
\end{tabular}}
\end{table}

\begin{table}[t]
\caption{Cascade Latency--Profit Frontier over the Boundary Band $\beta$ at
$\kappa=2$. Rollout share is the fraction of decisions escalated to the
rollout teacher.}
\label{tab:frontier}
\centering\small
\begin{tabular}{lrrrr}
\toprule
Scenario & $\beta$ & Retention & Mean latency ms & Rollout share \\
\midrule
\texttt{tight}  &   0 & 97.7\% &  6.27 & 28.6\% \\
\texttt{tight}  & 500 & 98.2\% & 12.95 & 44.7\% \\
\texttt{tight}  & 900 & 99.0\% & 17.72 & 57.0\% \\
\texttt{scarce} &   0 & 93.7\% &  6.09 & 45.4\% \\
\texttt{scarce} & 500 & 98.0\% & 11.54 & 59.6\% \\
\texttt{scarce} & 900 & 99.5\% & 15.72 & 78.6\% \\
\bottomrule
\end{tabular}
\end{table}

\begin{table}[t]
\caption{Full-Horizon Upper Bounds and Rollout Retention against Each
Ceiling. The Lagrangian-per-truck bound retains per-truck structure and is
20.7\% tighter than the LP relaxation on \texttt{tight}, 39.3\% on
\texttt{scarce}.}
\label{tab:bound}
\centering\small
\begin{tabular}{lrrrrr}
\toprule
Scenario & LP bound & Roll. ret. (LP) & Lagrangian bound & Roll. ret. (Lag) & Tightening \\
\midrule
\texttt{tight}  & \$2{,}377{,}500 & 53.6\% & \$1{,}885{,}043 & 67.6\% & 20.7\% \\
\texttt{scarce} & \$2{,}675{,}525 & 39.8\% & \$1{,}623{,}084 & 65.7\% & 39.3\% \\
\bottomrule
\end{tabular}
\end{table}

Three findings stand out. First, the cascade is the only dependency-free
policy that closes the gap to rollout: at $\beta=\$500$ it recovers about
98\% of rollout profit on both stress scenarios at 40\% (\texttt{tight}) to
56\% (\texttt{scarce}) of rollout's mean decision latency, and on
\texttt{tight} the paired cascade$-$rollout difference is not statistically
significant. Second, the frontier is smooth and monotone in $\beta$
(Table~\ref{tab:frontier}): the scarcity trigger alone recovers 97.7\% on
\texttt{tight} but only 93.7\% on \texttt{scarce}, where the surrogate's
disagreement with rollout extends beyond the scarcity regime. Third, the
standalone surrogate clears the best-simple bar on both scenarios but still
leaves 6--11 points of rollout value that only the cascade recovers: a
cheap model deployed alone under-performs precisely on the high-stakes
capacity decisions. The Lagrangian ceiling (Table~\ref{tab:bound})
relocates the rollout teacher from 53.6\%/39.8\% of the loose LP bound to
67.6\%/65.7\% of a structure-respecting bound; the residual gap is
dominated by inter-truck coordination slack, indicating that further
methods effort should target joint-decision lookahead rather than refining
single-tender scoring.

\section{Conclusion}
FreightBidBench provides a public, reproducible test bed for real-time
truckload bid acceptance under operational feasibility, with hindsight
ceilings that make achievable headroom explicit and a latency-aware cascade
that demonstrates a non-trivial latency--profit frontier. The benchmark,
ceilings, and versioned contracts let future methods---stronger surrogates,
joint-decision lookahead, learned escalation---be compared without
ambiguity over what is being measured. Extended proofs, full ablations,
sensitivity analyses, and calibration validation appear in the full
version.

\section*{Acknowledgments}
The author used Anthropic's Claude (Claude Opus), a large language model,
to assist with software implementation and experiment scaffolding,
drafting and editing of the manuscript, and literature and formatting
support. All model designs, experiments, results, and claims were verified
and approved by the author, who takes full responsibility for the content.

\section*{Author Contributions}
The author confirms sole responsibility for the following: study conception
and design; implementation; data curation; analysis and interpretation of
results; and draft manuscript preparation. The author reviewed the results
and approved the final version of the manuscript.

\section*{Declaration of Conflicting Interests}
The author declared the following potential conflict of interest with
respect to the research, authorship, and/or publication of this article:
the author is employed by Bubba AI, which develops AI-based load-planning
and carrier-operations products in the freight domain addressed by this
paper. This study uses only public FAF and USDA data and a
dependency-free, open-source benchmark; no proprietary data or systems were
used, and Bubba AI had no role in the study design, analysis, or the
decision to publish.

\section*{Funding}
The author disclosed no financial support for the research, authorship,
and/or publication of this article.

\section*{Data and Code Availability}
Code, data-processing pipelines, and run manifests are available at
\url{https://github.com/aswincsekar/freightbidbench}.

\bibliographystyle{apalike}
\bibliography{references}

\end{document}
